\section{Retained synthetic checks of signs and conditional identification}
\label{sec:synthetic}
\subsection{A local co-energy with physically admissible differential gains}
Use $x=i_d/I_*$, $y=i_q/I_*$, $\normalizedspeed=\omegae/\omega_*$, and
$\coenergy=\psi_*I_*\normalizedcoenergy$. Choose $I_*=\SI{100}{\ampere}$,
$\psi_*=\SI{0.1}{\weber}$, $\omega_*=\SI{1000}{\per\second}$, so
$L_*=\psi_*/I_*=\SI{1}{\milli\henry}$ and
$R_*=\omega_*\psi_*/I_*=\SI{1}{\ohm}$. The local polynomial is
\begin{equation}
 \normalizedcoenergy=x+\frac{0.6}{2}x^2+\frac12y^2
       -\frac{0.04}{4}x^4-\frac{0.08}{4}y^4
       -\frac{0.06}{2}x^2y^2+\frac{0.04}{2}xy^2.
 \label{eq:synthetic-energy}
\end{equation}
The linear term supplies PM flux. Negative fourth-order terms reduce d/q
differential gains; the two mixed terms introduce reciprocal cross-saturation
and an asymmetric d-bias dependence while retaining exact q reflection.
The resulting fluxes are
\begin{align}
 \psid/\psi_*&=1+0.6x-0.04x^3-0.06xy^2+0.02y^2,\\
 \psiq/\psi_*&=y-0.08y^3-0.06x^2y+0.04xy,
 \label{eq:synthetic-flux}
\end{align}
and the dimensionless Hessian is
\begin{equation}
 \Ldiff/L_*=
 \begin{bmatrix}
 0.6-0.12x^2-0.06y^2&-0.12xy+0.04y\\
 -0.12xy+0.04y&1-0.24y^2-0.06x^2+0.04x
 \end{bmatrix}.\label{eq:synthetic-hessian}
\end{equation}
Restrict this example to $|x|,|y|\leq1.1$. There the diagonal entries are
at least $0.3822$ and $0.5930$, while the off-diagonal magnitude is at most
$0.1892$. Strict diagonal dominance with positive diagonal proves positive
definiteness throughout that domain; the smaller Gershgorin bound is $0.1930$.
The polynomial is not a global high-current constitutive law: its negative
quartic terms would eventually destroy positive differential gains.

\subsection{Inject errors into voltage measurements, not into guessed curves}
For every observed current label, rotate to true current, compute true flux
and voltage with $R_s=\SI{0.08}{\ohm}$, rotate voltage back, and reconstruct
using the used resistance. Inject no error, $\resistanceerror=\SI{0.025}{\ohm}$ only,
$\angleerror=0.02$ rad only, or both. These steps reproduce
\cref{eq:combined-exact} by the measurement route independently of its map
derivation. The comparison in \cref{fig:synthetic} is at $x=0$, $\normalizedspeed=1$.
\begin{figure}[tbp]
\centering
\begin{tikzpicture}
\begin{groupplot}[group style={group size=2 by 2,horizontal sep=14mm,
 vertical sep=24mm},diagnostic axis,xmin=-1,xmax=1,
 legend style={compact text,at={(.5,-.32)},anchor=north,legend columns=4}]
\nextgroupplot[title={Reconstructed d flux},ylabel={$\hat\psi_d/\psi_*$}]
 \addplot[reference quantity,strong line] table[x=y,y=d0]{figures/data/synthetic.dat};
 \addlegendentry{Ideal}
 \addplot[estimated quantity,standard line,alternative pattern] table[x=y,y=dR]{figures/data/synthetic.dat};
 \addlegendentry{$R$}
 \addplot[derived quantity,standard line,sampled pattern] table[x=y,y=dg]{figures/data/synthetic.dat};
 \addlegendentry{Angle}
 \addplot[torque quantity,standard line,combined pattern] table[x=y,y=db]{figures/data/synthetic.dat};
 \addlegendentry{Both}
\nextgroupplot[title={Reconstructed q flux},ylabel={$\hat\psi_q/\psi_*$}]
 \addplot[reference quantity,strong line] table[x=y,y=q0]{figures/data/synthetic.dat};
 \addplot[estimated quantity,standard line,alternative pattern] table[x=y,y=qR]{figures/data/synthetic.dat};
 \addplot[derived quantity,standard line,sampled pattern] table[x=y,y=qg]{figures/data/synthetic.dat};
 \addplot[torque quantity,standard line,combined pattern] table[x=y,y=qb]{figures/data/synthetic.dat};
\nextgroupplot[title={Forbidden d-odd channel},ylabel={$r_d/\psi_*$}]
 \addplot[reference quantity,strong line] table[x=y,y=rd0]{figures/data/synthetic.dat};
 \addplot[estimated quantity,standard line,alternative pattern] table[x=y,y=rdR]{figures/data/synthetic.dat};
 \addplot[derived quantity,standard line,sampled pattern] table[x=y,y=rdg]{figures/data/synthetic.dat};
 \addplot[torque quantity,standard line,combined pattern] table[x=y,y=rdb]{figures/data/synthetic.dat};
\nextgroupplot[title={Forbidden q-even channel},ylabel={$r_q/\psi_*$},
 ymin=-.024,ymax=.003]
 \addplot[reference quantity,strong line] table[x=y,y=rq0]{figures/data/synthetic.dat};
 \addplot[estimated quantity,standard line,alternative pattern] table[x=y,y=rqR]{figures/data/synthetic.dat};
 \addplot[derived quantity,standard line,sampled pattern] table[x=y,y=rqg]{figures/data/synthetic.dat};
 \addplot[torque quantity,standard line,combined pattern] table[x=y,y=rqb]{figures/data/synthetic.dat};
\end{groupplot}
\end{tikzpicture}
\caption{Exact voltage-generated synthetic comparison at $x=0$, $\normalizedspeed=1$.
Injections are $\resistanceerror/R_*=0.025$ and $\angleerror=0.02$ rad. Parity channels
expose small errors that full flux curves partly hide. Resistance-only q flux
coincides with ideal q flux; the combined q-even residual coincides with
angle-only residual because resistance contributes zero q error at this label.}
\label{fig:synthetic}
\end{figure}


At $x=0$, the angle sensitivity in these normalized units is especially
instructive:
\begin{equation}
 \vect s_\gamma/\psi_*=
 \begin{bmatrix}0.4y-0.02y^3\\-1-0.06y^2\end{bmatrix}.
\end{equation}
The ideal d flux is $1+0.02y^2$, yet its odd residual remains zero.
The resistance tilt is $-0.025y/\normalizedspeed$. The angle changes both forbidden
channels: its d signal is slightly curved and its q signal grows with
$y^2$. Saturation here strengthens the q sensitivity and weakens the d
sensitivity relative to their small-current limits. It is not a universal
amplifier. The combined curves are additive in the resistance contribution
but nonlinear in angle.

For joint estimation use $x\in\{-0.6,0,0.6\}$, ten positive $y$ values
from $0.1$ to $1$, and $\normalizedspeed\in\{0.5,1,2\}$: 90 independent pairs and 180
residual entries. Fit the first-order columns using the known ideal synthetic
map. The noiseless estimates are shown below; they demonstrate sign recovery
and the small linearization bias, not uncertainty performance on real data.
\begin{table}[tb]
\centering\small\begin{tabular}{lrrr}
\toprule
Injected case & $\hat\varepsilon_R$ (\si{\ohm}) & $\widehat{\delta\gamma}$ (rad) & RMS$/\psi_*$\\
\midrule
No error & 0.000000 & 0.000000 & 2.51e-17 \\
Resistance only & 0.025000 & 0.000000 & 4.82e-17 \\
Angle only & 0.000000 & 0.019998 & 8.03e-07 \\
Both errors & 0.025000 & 0.019998 & 8.03e-07 \\
\bottomrule
\end{tabular}

\caption{First-order estimation from exact synthetic voltage reconstruction.
The nonzero injected values are $0.025$ ohm and $0.02$ rad. RMS is the
post-fit residual per scalar entry, normalized by $\psi_*$.}
\label{tab:estimates}
\end{table}

\subsection{Validation and the useful range of linearization}
The accompanying \texttt{numerics/validate.py} differentiates the co-energy
symbolically, checks reflection and reciprocity, differentiates the exact
rotation at zero angle, and checks all six reconstruction-input sensitivities.
It verifies zero-error recovery, resistance correction sign, isotropic and
zero-d-current limits, no-saturation and no-cross-saturation limits, and the
current-proportional voltage confounder. Numerical central differences check
the angle sensitivity independently over a $41\times81$ grid on $[-1,1]^2$.
All generated tables and plot data come from this script.

\begin{table}[tb]
\centering\small\begin{tabular}{rrr}
\toprule
$|\delta\gamma|$ (rad) & Max. map error$/\psi_*$ & Max. residual error$/\psi_*$\\
\midrule
0.001 & 8.974e-07 & 3.916e-10 \\
0.005 & 2.244e-05 & 4.895e-08 \\
0.010 & 8.975e-05 & 3.916e-07 \\
0.020 & 3.591e-04 & 3.133e-06 \\
0.050 & 2.245e-03 & 4.893e-05 \\
0.100 & 8.982e-03 & 3.911e-04 \\
\bottomrule
\end{tabular}

\caption{Maximum Euclidean discrepancy from first-order angle prediction
over the $41\times81$ grid. Exact rotated arguments stay in the stated
polynomial domain. These are absolute normalized errors, not relative errors
at a possibly vanishing signal.}
\label{tab:linearization}
\end{table}
The full map discrepancy in \cref{tab:linearization} scales quadratically
near zero. Its forbidden parity residual discrepancy scales cubically in this
smooth symmetric example: reflecting current relates the $+\angleerror$ and
$-\angleerror$ maps, so the forbidden channels are odd in angle and eliminate
even angle powers. This improvement is specific to exact symmetric pairing;
it does not excuse using a finite-angle linear model without checking it.
The script records numerical metrics in \texttt{numerics/validation.json}.

As an exact non-identifiability test, repeat the resistance-only case with
correct resistance and voltage error $\vect\varepsilon_u=-0.025R_*\current$.
It produces the same residuals to floating-point precision at every tested
speed. The least-squares result would look convincing but estimate an
effective voltage-drop mismatch rather than uniquely identify winding resistance.
